\chapter{空间解析几何}
\section{向量与坐标系}
\subsection{向量的定义与性质}
\begin{definition}
    空间中既有大小又有方向的量称为\overtext{向量}{vector}，记作$\vb{v}$.它的大小记为$|\vb{v}|$.
    \begin{itemize}
        \item 对于两个大小相等、方向相同的向量$\vb{a},\vb{b}$，称二者相等，记为$\vb{a} = \vb{b}$.
        \item 大小为$0$的向量称为\overtext{零向量}{zero vector}，记作$\vb{0}$，它没有确定的方向.
        \item 大小相等而方向相反的向量称为\overtext{负向量}{negative vector}，记作$-\vb{v}$.
        \item 大小为$1$的向量称为\overtext{单位向量}{unit vector}，记为$\vb{e_v} = \dfrac{\vb{v}}{|\vb{v}|}$.
    \end{itemize}
\end{definition}

\begin{enumerate}[property]
    \item 向量具有加法且具备如下性质：
    \begin{itemize}
        \item \textbf{交换律：}$\vb{a + b = b + a}$.
        \item \textbf{结合律：}$\vb{(a + b) + c = a + (b + c)}$.
        \item \textbf{零元与逆元：}$\vb{a + 0 = a},\quad \vb{a + (-a) = 0}$.
    \end{itemize}
    上述性质利用平行四边形法则或三角形法则易证.
    \item 向量具有数乘.设$\lambda,\mu \in \mathbb{R}$，满足：
    \begin{itemize}
        \item $\lambda(\vb{a+b}) = \lambda\vb{a} + \lambda\vb{b}$.
        \item $(\lambda + \mu)\vb{a} = \lambda\vb{a} + \mu\vb{a}$.
        \item $(\lambda\mu)\vb{a} = \lambda(\mu\vb{a})$.
    \end{itemize}
\end{enumerate}

\begin{theorem}
    设$\vb{a,b}$是两个向量，则：
    \[\vb{|a+b| \leqslant |a| + |b|}\]
    当且仅当存在非负实数$\lambda$满足$\vb{a} = \lambda\vb{b}$，等号成立.
\end{theorem}


\subsection{向量的共线与共面}
\begin{definition}
    如果通过平移：
    \begin{itemize}
        \item 能使一组向量$\vb{a,b}$同处一条直线上，那么称它们是\textbf{共线}或\overtext{平行}{parallel}的，记为$\vb{a}\parallel\vb{b}$.
        \item 能使一组向量$\vb{a,b,c}$同处一个平面上，那么称它们是\textbf{共面}的.
    \end{itemize}
\end{definition}

\begin{theorem}
    \begin{itemize}
        \item 两个向量$\vb{a,b}$共线，当且仅当它们线性相关，即存在不全为零的实数$\lambda,\mu$满足：
        \[\lambda\vb{a} + \mu\vb{b} = \vb{0}\]
        \item 三个向量$\vb{a,b,c}$共面，当且仅当它们线性相关，即存在三个不全为零的实数$\lambda,\mu,\nu$满足：
        \[\lambda\vb{a} + \mu\vb{b} + \nu\vb{c} = \vb{0}\]
    \end{itemize}
\end{theorem}

\begin{proof}
    以两个向量$\vb{a,b}$的证明为例：

    \textbf{充分性：}设非零向量$\vb{a,b}$共线，即：
    \[\vb{a} = \pm k \vb{b},\quad k = \frac{|\vb{a}|}{|\vb{b}|} \neq 0 \quad\Rightarrow\quad \vb{a} \mp k \vb{b} = \vb{0}\]
    取$\lambda = 1 \neq 0,\mu = \mp k \neq 0$，满足：
    \[\lambda \vb{a} + \mu \vb{b} = \vb{0}\]

    \textbf{必要性：}若存在不全为零的数$\lambda,\mu$使得：
    \[\lambda \vb{a} + \mu \vb{b} = \vb{0}\]
    不妨设$\lambda \neq 0$，则有：
    \[\vb{a} = \frac{\mu}{\lambda}\vb{b}\]
    满足共线的定义.

    对于多个向量的证明，只需要对向量进行组合，构造共线向量即可.
\end{proof}


\subsection{向量的内积和外积}
\subsubsection{内积}
设$\vb{a,b}$是两个向量，它们的夹角定义为$\theta = \theta(\vb{a,b}) \in [0,\pi]$，显然有$\theta(\vb{a,b}) = \theta(\vb{b,a})$.对于任意非零的实数$\lambda$满足：
\[\theta(\lambda\vb{a,b}) = \begin{cases}
    \theta(\vb{a,b}),\quad\lambda>0\\
    \pi-\theta(\vb{a,b}),\quad\lambda<0
\end{cases}\]

\begin{definition}
    定义两个向量$\vb{a,b}$的\overtext{内积}{inner product}\!\!\textbf{(\!\!\overtext{点乘}{dot product}\!\!、\overtext{数量积}{scalar product}\!\!)}为如下运算：
    \[\vb{a}\cdot\vb{b} = (\vb{a},\vb{b}) \coloneq |\vb{a}| |\vb{b}| \cos\theta(\vb{a,b}) \in \mathbb{R},\quad \theta \in [0,\pi]\]
    若其中一个向量是零向量，则规定点乘结果为$0$.
\end{definition}

显然，从定义上看，总是有：
\[\begin{cases}
    \vb{a \cdot b} > 0,\quad \theta(\vb{a,b}) < \dfrac{\pi}{2} \\[0.5em]
    \vb{a \cdot b} = 0,\quad \theta(\vb{a,b}) = \dfrac{\pi}{2} \quad\Rightarrow\quad \text{此时称$\vb{a,b}$\overtext{正交}{orthogonal}或\overtext{垂直}{vertical}，记作$\vb{a}\perp\vb{b}$}
\end{cases}\]
几何上，$|\vb{a}|\cos\theta$就是$\vb{a}$在$\vb{b}$所在直线上投影向量$\vb{a_b}$的有向长度，有如下关系式：
\[\vb{a_b} = (|\vb{a}|\cos\theta)\vb{e_b} = (\vb{a \cdot e_b})\vb{e_b}\]
上式表明当$\theta \in [\frac{\pi}{2},\pi]$时，投影向量$\vb{e_b}$与$\vb{b}$共线但反向.

于是点乘的含义可以理解为$\vb{a}$在$\vb{b}$方向上的投影向量的有向长度乘$\vb{b}$的长度，或是说与$\vb{b}$垂直的分量对$\vb{a \cdot b}$没有贡献.

点乘具有如下性质：
\begin{itemize}
    \item \textbf{交换律：}$\vb{a \cdot b = b \cdot a}$.
    \item \textbf{结合律：}$(\lambda \vb{a})\cdot \vb{b} = \lambda (\vb{a \cdot b})$.
    \item \textbf{分配律：}$\vb{(a + b)\cdot c = a \cdot c + b \cdot c}$
\end{itemize}


\subsubsection{外积}
对于不共面向量组$\{\vb{a,b,c}\}$，若其满足右手系，则$\{\vb{b,c,a}\},\{\vb{c,a,b}\}$仍是右手系，但$\{\vb{b,a,c}\}$是左手系.

\begin{definition}
    定义两个向量$\vb{a,b}$的\overtext{外积}{outer product}\!\!\textbf{(\!\!\overtext{叉乘}{cross product}\!\!、\overtext{向量积}{vector product}\!\!)}的结果为一个矢量，其大小规定为：
    \[|\vb{a \times b}| \coloneq |\vb{a}||\vb{b}|\sin\theta(\vb{a,b}),\quad \theta\in [0,\pi]\]
    方向规定为$\vb{a \times b}$与$\vb{a,b}$垂直，且$\{\vb{a,b,a \times b}\}$构成右手系.

    若其中一个向量是零向量，则规定叉乘结果也是零向量.
\end{definition}

显然，从定义上看，若是两个向量贡献，其叉乘结果是零向量.

从几何上看，外积$\vb{a \times b}$的大小即为以$\vb{a,b}$为邻边的平行四边形的面积.

叉乘具有如下性质：
\begin{itemize}
    \item $(\lambda\vb{a})\times\vb{b} = \vb{a}\times(\lambda\vb{b}) = \lambda(\vb{a}\times\vb{b})$.
    \item \textbf{反称性：}$\vb{a \times b = -a \times b}$.
    \item \textbf{分配律：}$\vb{a \times (b + c) = a \times b + a \times c}$.
\end{itemize}


\subsubsection{混合积}
\begin{definition}
    称$\vb{(a \times b)\cdot c}$为三个向量$\vb{a,b,c}$的\overtext{混合积}{mixed product}.
\end{definition}

从几何上看，混合积结果的绝对值是以$\vb{a,b,c}$为棱的平行六面体的体积.

若是有$\vb{(a \times b)\cdot c} = 0$，则表明$\vb{a,b,c}$共面.

混合积具有\overtext{轮换对称性}{cyclic symmetry}，即：
\[\vb{a \times b}\cdot c = \vb{b \times c}\cdot a = \vb{c \times a}\cdot b\]


\subsubsection{有向长度、面积和体积}
\begin{definition}
    有向长度、面积和体积即为具备正负性的长度、面积和体积，规定：
    \begin{itemize}
        \item 当$\theta(\vb{a,b})$为锐角时，有向长度$|\vb{a}|\cos\theta$为正；当$\theta(\vb{a,b})$为钝角时为负.
        \item 当$\{\vb{a,b,a \times b}\}$构成右手系时，有向面积$\vb{a \times b}$为正，反之构成左手系则为负.
        \item 当$\{\vb{a,b,c}\}$构成右手系时，有向体积$\vb{(a \times b)\cdot c}$为正，反之构成左手系则为负.
    \end{itemize}
\end{definition}


\subsection{向量的坐标表示与空间坐标系}
\begin{definition}
    设集合$V$由全部空间向量构成的集合，在其上定义了向量的加法和数乘，称$V$是一个\overtext{线性空间}{linear space}\textbf{(向量空间)}.

    称$V$中三个不共面的向量组$\{\vb{e}_1,\vb{e}_2,\vb{e_3}\}$为$V$的一组\overtext{基向量}{basis vectors}，简称\overtext{基}{base}.

    选定一个原点$O$，称$\mathbb{F} = [O;\vb{e}_1,\vb{e}_2,\vb{e}_3]$为$V$中的一个\overtext{仿射坐标系}{affine coordinate system}，其中一向量$\vb{x}$对应的坐标$(x_1,x_2,x_3)$称为\overtext{仿射坐标}{affine coordinates}.
\end{definition}

\begin{theorem}
    若线性空间$V$中存在一组基$\{\vb{e}_1,\vb{e}_2,\vb{e}_3\}$，则对$\forall \vb{x} \in V$，都存在唯一的一组数$(x_1,x_2,x_3)$使得：
    \[\vb{x} = x_1\vb{e}_1 + x_2\vb{e}_2 + x_3\vb{e}_3 \quad\text{或}\quad \vb{x} \mapsto (x_1,x_2,x_3)\]
\end{theorem}

\begin{definition}
    取线性空间$V$中的$\{\vb{i,j,k}\}$是两两正交且形成右手系的单位向量，称他们为$V$的\overtext{标准正交基}{orthonormal basis}，称$\mathbb{F} = [O;\vb{i,j,k}]$为\overtext{直角坐标系}{cartesian coordinate system}.
\end{definition}

直角坐标系中满足如下等式：
\begin{gather*}
    \vb{i \cdot i} = \vb{j \cdot j} = \vb{k \cdot k} = 1 \\
    \vb{i \cdot j} = \vb{j \cdot k} = \vb{k \cdot i} = 0 \\
    \vb{i \times i} =\vb{j \times j} = \vb{k \times k} = \vb{0} \\
    \vb{i \times j} = \vb{k},\quad \vb{j \times k} = \vb{i},\quad \vb{k \times i} = \vb{j}
\end{gather*}
根据上述内容不难推出内积、外积和混合积的坐标表示：
\[\vb{a \cdot b} = a_1b_1 + a_2b_2 + a_3b_3,\quad \vb{a \times b} = \begin{vmatrix}
    \vb{i} & \vb{j} & \vb{k} \\
    a_1 & a_2 & a_3 \\
    b_1 & b_2 & b_3
\end{vmatrix},\quad \vb{(a \times b)\cdot c} = \begin{vmatrix}
    a_1 & a_2 & a_3 \\
    b_1 & b_2 & b_3 \\
    c_1 & c_2 & c_3
\end{vmatrix}\]

\begin{theorem}
    \textbf{Cauchy-Schwarz不等式}
    \[(a_1b_1 + a_2b_2 + a_3b_3)^2 \leqslant (a_1^2 + a_2^2 + a_3^2)(b_1^2 + b_2^2 + b_3^2)\]
\end{theorem}

\begin{proof}
    设有$\vb{a} = (a_1,a_2,a_3),\vb{b}=(b_1,b_2,b_3)$，则不等式等价为：
    \[(\vb{a \cdot b})^2 = (|\vb{a}||\vb{b}|\cos\theta)^2 \leqslant |\vb{a}|^2 |\vb{b}|^2\]
    显然得证.
\end{proof}

\begin{definition}
    在空间中取定一点$O$，则空间中任意一点$P$对应一个向量$\va{OP}$，称为$P$的\overtext{位置向量}{positional vector}，记作$\vb{r}(P)$，在直角坐标系下具有如下对应关系：
    \[P \leftrightarrow \va{OP} = x\vb{i} + y\vb{j} + z\vb{k} \leftrightarrow (x,y,z)\]
    $\va{OP}$的方向可由$\va{OP}$与三个坐标轴的夹角，即\overtext{方向角}{direction angles}来刻画：
    \[x = |\va{OP}|\cos\alpha,\quad y = |\va{OP}|\cos\beta,\quad z = |\va{OP}|\cos\gamma\]
    称方向角的余弦为\overtext{方向余弦}{direction cosines}.
\end{definition}
注意到，$\va{OP}$的方向余弦构成了$\va{OP}$方向上的一个单位向量$(\cos\alpha,\cos\beta,\cos\gamma)$.


\section{坐标变换和其它常用坐标系}
\subsection{坐标变换}
\subsubsection{平移}
将坐标系$\mathbb{F}$平移得到的新坐标系$\mathbb{F'}$，平移向量为$\va{OO'} = a\vb{i} + b\vb{j} + c\vb{k}$.设$P$在两个坐标系下的位置向量为：
\[\va{OP} = x\vb{i} + y\vb{j} + z\vb{k},\qquad \va{O'P} = x'\va{i} + y'\va{j} + z'\va{k}\]
则根据$\va{O'P} = \va{OP} - \va{OO'}$可得平移后的坐标：
\[x'=x-a,\quad y'=y-b,\quad z'=z-c\]


\subsubsection{旋转}
假设$\mathbb{F'}$由$\mathbb{F}$保持原点不动旋转得到，设两个坐标系基向量两两之间的夹角为：
\begin{gather*}
    \begin{NiceArray}{c|ccc}
        & \vb{i} & \vb{j} & \vb{k} \\
        \Hline
        \vb{i'} & \alpha_1 & \beta_1 & \gamma_1 \\
        \vb{j'} & \alpha_2 & \beta_2 & \gamma_2 \\
        \vb{k'} & \alpha_3 & \beta_3 & \gamma_3 \\
    \end{NiceArray} \quad\Rightarrow\quad \begin{cases}
        \vb{i'} = \cos\alpha_1\vb{i} + \cos\beta_1\vb{j} + \cos\gamma_1\vb{k} \\
        \vb{j'} = \cos\alpha_2\vb{i} + \cos\beta_2\vb{j} + \cos\gamma_2\vb{k} \\
        \vb{k'} = \cos\alpha_3\vb{i} + \cos\beta_3\vb{j} + \cos\gamma_3\vb{k}
    \end{cases}\\
    \begin{bmatrix}
        \cos\alpha_1 & \cos\beta_1 & \cos\gamma_1 \\
        \cos\alpha_2 & \cos\beta_2 & \cos\gamma_2 \\
        \cos\alpha_3 & \cos\beta_3 & \cos\gamma_3
    \end{bmatrix} \begin{bmatrix}
        \vb{i} \\ \vb{j} \\ \vb{k}
    \end{bmatrix} =  \begin{bmatrix}
        \vb{i'} \\ \vb{j'} \\ \vb{k'}
    \end{bmatrix} 
\end{gather*}
也就是说，左边的$3 \times 3$矩阵实现了旋转变换操作.


\subsection{其它坐标系}
\begin{itemize}
    \item \textbf{极坐标：}极坐标$(r,\theta)$与直角坐标$(x,y)$之间的变换关系为：
    \[\begin{cases}
        x = r\cos\theta \\
        y = r\sin\theta
    \end{cases} \Leftrightarrow\quad \begin{cases}
        r = \sqrt{x^2 + y^2} \\
        \theta = \arctan\dfrac{y}{x}
    \end{cases}\]
    \item \textbf{柱坐标：}柱坐标$(r,\theta,z)$与直角坐标$(x,y,z)$之间的变换关系为：
    \[x = r\cos\theta,\quad y = r\sin\theta,\quad z = z\]
    \item \textbf{球坐标：}球坐标$(r,\theta,\varphi)$与直角坐标$(x,y,z)$之间的变换关系为：
    \[x = r\sin\theta\cos\varphi,\quad y = r\sin\theta\sin\varphi,\quad z = r\cos\theta\]
\end{itemize}


\section{平面与直线}
\subsection{平面方程}
\subsubsection{平面的一般方程}
平面的一般方程来自于如下方式的构造：在给定一点$P_0(x_0,y_0,z_0)$上延伸出一个向量$\vb{n} = (a,b,c)$，取一个平面过点$P_0$且垂直于这个向量.

按照上述构造的过程，取平面上任意一点$P(x,y,z)$，必然有$\vb{n}\cdot\va{P_0P} = 0$，即：
\[a(x-x_0) + b(y-y_0) + c(z-z_0) = 0\]
令$d \equiv -(ax_0 + by_0 + cz_0)$，即：
\[ax + by + cz + d = 0\]
满足上式的点$(x,y,z)$一定在上述平面上，上式称为\textbf{平面的一般方程}.

上式也意味着，对于表示为平面的一般方程的平面，取$x,y,z$的系数即位其法向量.


\subsubsection{平面的截距式方程}
平面的截距式方程来自于如下方式的构造：对于不共线的三个点$P_1,P_2,P_3$，必有$\va{P_1P_2},\va{P_1P_3}$形成一个平面，所以由外积可以得到这个平面点法向量$\va{P_1P_2}\times\va{P_1P_3}$.

取平面上任意一点$P(x,y,z)$，总是有：
\[(\va{P_1P_2} \times \va{P_1P_3}) \cdot \va{P_1P} = \begin{vmatrix}
    x_2 - x_1 & y_2-y_1 & z_2-z_1 \\
    x_3 - x_1 & y_3-y_1 & z_3-z_1 \\
    x - x_1 & y-y_1 & z-z_1
\end{vmatrix} = 0\]
特别地，若是取三个坐标轴上的点$(\alpha,0,0),(0,\beta,0),(0,0,\gamma)$构造平面，上述方程简化为\textbf{平面的截距式方程}：
\[\frac{x}{\alpha} + \frac{y}{\beta} + \frac{z}{\gamma} = 1\]


\subsubsection{两平面的关系}
平面间相交与平行的关系主要是通过法向量之间的关系来确定的：
\begin{itemize}
    \item \textbf{平行：}若两平面平行，则有$\vb{n}_1 \parallel \vb{n}_2$.即存在$\lambda \neq 0$使得$\vb{n}_1 = \lambda \vb{n}_2$.
    
    对于平面的一般方程，有$a_1 = \lambda a_2, b_1 = \lambda b_2, c_1 = \lambda c_2$，若同时又有$d_1 = \lambda d_2$，则两平面重合.
    \item \textbf{相交：}两平面的法向量不平行，则两平面相交，相交的二面角为：
    \[\theta = \arccos\frac{|\vb{n}_1 \cdot \vb{n}_2|}{|\vb{n}_1||\vb{n}_2|}\]
\end{itemize}


\subsubsection{点到平面的距离}
设有平面外一点$P_0(x_0,y_0,z_0)$，过点$P_0$作平面的垂线，垂足为$P_1$.在平面上任取一点$P(x,y,z)$，设$\va{PP_0}$与平面法向量$\vb{n} = (a,b,c)$的夹角为$\theta$，则有：
\begin{align*}
    |\va{P_0P_1}| &= |\va{PP_0}||\cos\theta| = \left|\frac{\vb{n}}{|\vb{n}|}\cdot\va{PP_0}\right| = \frac{|a(x_0-x) + b(y_0-y) + c(z_0-z)|}{\sqrt{a^2 + b^2 + c^2}}\\
    &= \frac{|ax_0 + by_0 + cz_0 +d|}{\sqrt{a^2 + b^2 + c^2}}
\end{align*}
其分子相当于把平面外一点的坐标代入平面方程.

空间中的一张平面将空间分为平面法向所指的一侧(上方)、平面本身、法向所指的另一侧(下方).判断一个点属于哪个部分，主要是看$\va{PP_0}$与法向量$\vb{n}$的夹角：
\[\vb{n}\cdot\va{PP_0} = ax_0 + by_0 + cz_0 + d\begin{cases}
    >0 \quad\Rightarrow\quad \text{是锐角，$P_0$在平面的上方.}\\
    =0 \quad\Rightarrow\quad \text{是直角，$P_0$在平面内.}\\
    <0 \quad\Rightarrow\quad \text{是钝角，$P_0$在平面的下方.}
\end{cases}\]


\subsection{直线方程}
\subsubsection{直线的点向式方程}
直线的点向式方程来自于如下方式的构造：设直线$L$过一点$P_0(x_0,y_0,z_0)$且其平行于方向向量$\vb{v} = (l,m,n)$，则对于平面上任意一点$P(x,y,z)$有：
\[\va{OP_0} - \va{OP} = t\vb{v} \quad\Rightarrow\quad \begin{cases}
    x = x_0 + lt\\
    y = y_0 + mt\\
    z = z_0 + nt
\end{cases}\]
上述方程组为直线的\textbf{参数方程}，将其联立消去$t$则得到直线的\textbf{点向式方程}：
\[\frac{x-x_0}{l} = \frac{y-y_0}{m} = \frac{z-z_0}{n}\]
根据上述方程构造的原理，上式中若有$l,m,n$为$0$的情况，应等价于对应分子等于$0$.

设若直线过给定点$P_1(x_1,y_1,z_1),P_2(x_2,y_2,z_2)$，其点向式方程为：
\[\frac{x-x_1}{x_2-x_1} = \frac{y-y_1}{y_2-y_1} = \frac{z-z_1}{z_2-z_1}\]


\subsubsection{直线的一般方程}
直线的一般方程来自于如下方式的构造：联立两个平面的一般方程，则得到它们的交线.设交线上一点$P(x,y,z)$，其应满足：
\[\begin{cases}
    a_1x + b_1y + c_1z + d_1 = 0\\
    a_2x + b_2y + c_2z + d_2 = 0
\end{cases}\]
上式称为\textbf{直线的一般方程}，其方向向量$\vb{v} = \vb{n}_1 \times \vb{n}_2$.




\subsubsection{两直线的位置关系}
对于给定的两条直线，设其上各有任意一点$P_1(x_1,y_1,z_1),P_2(x_2,y_2,z_2)$，其共面的条件是$\va{P_1P_2},\vb{v}_1,\vb{v}_2$三向量共面，即其混合积为$0$：
\[(\vb{v}_1 \times \vb{v}_2) \cdot \va{P_1P_2} = \begin{vmatrix}
    x_2 - x_1 & y_2-y_1 & z_2-z_1 \\
    l_1 & m_1 & n_1 \\
    l_2 & m_2 & n_2
\end{vmatrix} = 0\]
则两直线不共面的条件则为上式不为$0$.

对于共面的两直线，其夹角$\theta$满足：
\[\cos\theta = \frac{\vb{v}_1 \cdot \vb{v}_2}{|\vb{v}_1||\vb{v}_2|} \begin{cases}
    = \pm 1 \quad\Rightarrow\quad \text{两直线平行.}\\
    \in (-1,1) \quad\Rightarrow\quad \text{两直线相交.}
\end{cases}\]


\subsubsection{点到直线的距离}
设有直线外一点$P_0(x_0,y_0,z_0)$、直线上任意一点$P(x,y,z)$，则$P_0$到直线的距离为：
\[d = |\va{PP_0}|\sin\theta = \frac{|\vb{v}\times\va{PP_0}|}{|\vb{v}|}\]


\subsubsection{直线与平面的位置关系}
直线与平面的关系主要是通过直线与平面的交角来确定的.设平面法向量$\vb{n}$与直线方向向量$\vb{v}$的夹角为$\theta$，则直线与平面的交角$\varphi = \frac{\pi}{2}-\theta$或$-\frac{\pi}{2}+\theta$，于是有：
\[\sin\varphi = |\cos\theta| = \frac{|\vb{n \cdot v}|}{|\vb{n}||\vb{v}|} \begin{cases}
    = 1 \quad\Rightarrow\quad \text{直线垂直于平面.}\\
    \in (0,1) \quad\Rightarrow\quad \text{直线与平面相交.}\\
    = 0 \quad\Rightarrow\quad \text{直线平行于平面.}\\
\end{cases}\]


\section{二次曲面}
\subsection{柱面}
\begin{definition}
    曲线$F(x,y) = 0$在空间坐标系下$Oxy$平面内的表达式为：
    \[\begin{cases}
        F(x,y) = 0\\
        z = 0
    \end{cases}\]
    上述曲线称为\overtext{准线}{directrix}.取一条垂直于$Oxy$平面的直线，使其沿着准线移动，形成的面称为\overtext{柱面}{cylindrical surface}，其方程为：
    \[F(x,y) = 0\]
    称动垂线为这个柱面的\overtext{母线}{generatrix}.
\end{definition}

常见的圆锥曲线方程在空间坐标下则变成相应的柱面：
\begin{itemize}
    \item \textbf{椭圆柱面：}$\dfrac{x^2}{a^2} + \dfrac{y^2}{b^2} = 1$.
    \item \textbf{双曲柱面：}$\dfrac{x^2}{a^2} - \dfrac{y^2}{b^2} = 1$.
    \item \textbf{抛物柱面：}$y^2 - 2px = 0,\quad p>0$.
\end{itemize}

\begin{figure}[H]
    \centering
    \includegraphics[width=\textwidth]{quadratic-surface/cylinderical-surface/fig.pdf}
    \caption{椭圆柱面、双曲柱面与抛物柱面及其准线}
\end{figure}


\subsection{旋转面}
\begin{definition}
    设空间中$Oyz$平面上有一曲线$\begin{cases}
        F(y,z) = 0,\quad y>0\\
        x = 0
    \end{cases}$，使其绕$z$轴旋转得到一个曲面，称为\overtext{旋转曲面}{rotating surface}，其方程为：
    \[F(\pm\sqrt{x^2 + y^2},z) = 0\]
    即绕哪个轴旋转，就将平面曲线方程的另一个变量替换为圆的方程.
\end{definition}

\begin{figure}[H]
    \centering
    \includegraphics[width=\textwidth]{quadratic-surface/rotating-surface/fig.pdf}
    \vspace{-3em}
    \caption{旋转曲面}
\end{figure}


\subsection{椭球面}
\begin{definition}
    由下述方程描述的曲面称为\overtext{椭球面}{ellipsoid surface}：
    \[\frac{x^2}{a^2} + \frac{y^2}{b^2} + \frac{z^2}{c^2} = 1,\quad a,b,c>0\]
\end{definition}

一般的椭球面并不一定是旋转面，只有三个半轴中有两个一样长的才是旋转椭球面，若有三个一样长的半轴，则退化为球面.

\begin{figure}[H]
    \centering
    \includegraphics[width=\textwidth]{quadratic-surface/ellipsoid-surface/fig.pdf}
    \vspace{-3em}
    \caption{椭球面、旋转椭球面和球面}
\end{figure}


\subsection{双曲面}
\begin{definition}
    由下述方程描述的曲面称为\overtext{单叶双曲面}{hyperboloid of one sheet}： 
    \[\frac{x^2}{a^2} + \frac{y^2}{b^2} - \frac{z^2}{c^2} = 1,\quad a,b,c>0\]
    含有两个减号，由下述方程描述的曲面称为\overtext{双叶双曲面}{hyperboloid of two sheets}： 
    \[\frac{x^2}{a^2} - \frac{y^2}{b^2} - \frac{z^2}{c^2} = 1,\quad a,b,c>0\]
    双叶双曲线上离原点最近的点称为双叶双曲面的\overtext{顶点}{vertex}
\end{definition}


\subsection{二次锥面}
\begin{definition}
    由下述方程描述的曲面称为\overtext{二次锥面}{quadric cone}：
    \[\frac{x^2}{a^2} + \frac{y^2}{b^2} - \frac{z^2}{c^2} = 0,\quad a,b,c>0\]
    二次锥面由通过原点$O$的与平面交角相同的直线所构成，称为锥面的\textbf{母线}，点$O$称为锥面的\textbf{顶点}.
\end{definition}

\begin{figure}[H]
    \centering
    \includegraphics[width=\textwidth]{quadratic-surface/hyperboloid-and-cone-surface/fig.pdf}
    \vspace{-2.25em}
    \caption{单叶双曲面、双叶双曲面和二次锥面}
\end{figure}


\subsection{椭圆抛物面}
\begin{definition}
    由下述方程定义的曲面称为\overtext{椭圆抛物面}{elliptic paraboloid}：
    \[\frac{x^2}{a^2} + \frac{y^2}{b^2} = z,\quad a,b>0\]
\end{definition}


\subsection{双曲抛物面(马鞍面)}
\begin{definition}
    由下述方程定义的曲面称为\overtext{双曲抛物面}{hyperbolic paraboloid}，也称作\overtext{马鞍面}{saddle surface}：
    \[\frac{x^2}{a^2} - \frac{y^2}{b^2} = z,\quad a,b>0\]
\end{definition}

\begin{figure}[H]
    \centering
    \includegraphics[width=\textwidth]{quadratic-surface/paraboloid-surface/fig.pdf}
    \vspace{-5.25em}
    \caption{椭圆抛物面和双曲抛物面}
\end{figure}

实际上，上述二次曲面都是一个一般的三元二次方程的退化或经坐标变换后的形式：
\[a_{11}x^2 + a_{22}y^2 + a_{33}z^2 + 2a_{12}xy + 2a_{13}xz + 2a_{23}yz + 2b_1x + 2b_2y + 2b_3z + c = 0\]